% GENERATED FILE. Edit canonical lessons, then run npm run generate:books.
% canonical-source-hash: fnv1a64:4816c61dd8264b52
% canonical-lessons: ML-C485-liphrru, ML-C485-konippati, ML-C485-kavatam, ML-C485-visramamuri, ML-C485-tattukata

\chapter{Ways In and Out}
\label{ch:ml-ways-in-and-out}
\hlchaptermodality{485}

\begin{chapteropening}
\textbf{By the end of this chapter:} \emph{I can say a lift, a staircase, a gateway, a waiting room and a roadside food stall.}

Complete the last lesson of chapter 485: I can say a lift, a staircase, a gateway, a waiting room and a roadside food stall.
\end{chapteropening}

\section[\texorpdfstring{\ml{ലിഫ്റ്റ്} (liphṟṟŭ)}{ലിഫ്റ്റ് (liphṟṟŭ)}]{\ml{ലിഫ്റ്റ്} (liphṟṟŭ) — a lift}

\label{lesson:ML-C485-liphrru}



\begin{quote}
Before the new one: say the Malayalam for a window, then the Malayalam for a view.
\end{quote}

\subsection*{You'll want to know: \ml{ലിഫ്റ്റ്}}
\textbf{\ml{ലിഫ്റ്റ്}} — \emph{liphṟṟŭ} — \textquotedblleft{}a lift\textquotedblright{}.

\begin{quote}
\ml{ലിഫ്റ്റ്} \ml{കേടായതുകൊണ്ട്} \ml{ഞങ്ങൾ} \ml{നടന്നു} \ml{കയറി}.
\end{quote}

\begin{quote}
\emph{liphṟṟŭ kēṭāyatukoṇṭŭ ñaṅṅaḷ naṭannu kayaṟi.}
\end{quote}

\begin{quote}
Because the lift had broken down, we walked up.
\end{quote}

British usage: Malayalam says \ml{ലിഫ്റ്റ്}, not elevator.

\begin{cousinweb}
\textbf{\ml{ലിഫ്റ്റ്}} (\emph{liphṟṟŭ}) was borrowed from English \emph{lift}, a moving box that carries people between floors.
\end{cousinweb}

\begin{grammarlens}[title={What you've built}]
One more word for getting around.
\end{grammarlens}

\subsection*{Your turn}
\begin{itemize}
\raggedright
  \item \emph{Say it:} \emph{liphṟṟŭ}
  \item \emph{Say it:} \emph{liphṟṟŭ}, once more
  \item \emph{Say it:} \emph{liphṟṟŭ kēṭāyatukoṇṭŭ ñaṅṅaḷ naṭannu kayaṟi.} — Because the lift had broken down, we walked up.
  \item \emph{Recall:} say the Malayalam for a window, then the Malayalam for a view, then say \emph{liphṟṟŭ} again
\end{itemize}

\begin{tcolorbox}[breakable,colback=teal!4,colframe=teal!35!black,title={Before you move on}]
What does \ml{ലിഫ്റ്റ്} mean? (\textquotedblleft{}A lift\textquotedblright{}.) Say it once more.
\end{tcolorbox}
\section[\texorpdfstring{\ml{കോണിപ്പടി} (kōṇippaṭi)}{കോണിപ്പടി (kōṇippaṭi)}]{\ml{കോണിപ്പടി} (kōṇippaṭi) — a staircase}

\label{lesson:ML-C485-konippati}



\begin{quote}
Before the new one: say the Malayalam for a view, then the Malayalam for a lift.
\end{quote}

\subsection*{You'll want to know: \ml{കോണിപ്പടി}}
\textbf{\ml{കോണിപ്പടി}} — \emph{kōṇippaṭi} — \textquotedblleft{}a staircase\textquotedblright{}.

\begin{quote}
\ml{ഞങ്ങൾ} \ml{കോണിപ്പടി} \ml{കയറി} \ml{മൂന്നാം} \ml{നിലയിലെത്തി}.
\end{quote}

\begin{quote}
\emph{ñaṅṅaḷ kōṇippaṭi kayaṟi mūnnāṁ nilayiletti.}
\end{quote}

\begin{quote}
We climbed the stairs and reached the third floor.
\end{quote}

\ml{കോണി}, a ladder, plus \ml{പടി}, a step. One step is a \ml{പടി}.

\begin{cousinweb}
\textbf{\ml{കോണിപ്പടി}} (\emph{kōṇippaṭi}) is a Malayalam compound built from \ml{കോണി}, a ladder, and the native word \ml{പടി}, a step or stair (DEDR 3850).
\end{cousinweb}

\begin{grammarlens}[title={What you've built}]
One more word for getting around.
\end{grammarlens}

\subsection*{Your turn}
\begin{itemize}
\raggedright
  \item \emph{Say it:} \emph{kōṇippaṭi}
  \item \emph{Say it:} \emph{kōṇippaṭi}, once more
  \item \emph{Say it:} \emph{ñaṅṅaḷ kōṇippaṭi kayaṟi mūnnāṁ nilayiletti.} — We climbed the stairs and reached the third floor.
  \item \emph{Recall:} say the Malayalam for a view, then the Malayalam for a lift, then say \emph{kōṇippaṭi} again
\end{itemize}

\begin{tcolorbox}[breakable,colback=teal!4,colframe=teal!35!black,title={Before you move on}]
What does \ml{കോണിപ്പടി} mean? (\textquotedblleft{}A staircase\textquotedblright{}.) Say it once more.
\end{tcolorbox}
\section[\texorpdfstring{\ml{കവാടം} (kavāṭaṁ)}{കവാടം (kavāṭaṁ)}]{\ml{കവാടം} (kavāṭaṁ) — a gateway; an entrance}

\label{lesson:ML-C485-kavatam}



\begin{quote}
Before the new one: say the Malayalam for a lift, then the Malayalam for a staircase.
\end{quote}

\subsection*{You'll want to know: \ml{കവാടം}}
\textbf{\ml{കവാടം}} — \emph{kavāṭaṁ} — \textquotedblleft{}a gateway; an entrance\textquotedblright{}.

\begin{quote}
\ml{സ്റ്റേഷന്റെ} \ml{പ്രധാന} \ml{കവാടത്തിൽ} \ml{നമുക്ക്} \ml{കാണാം}.
\end{quote}

\begin{quote}
\emph{sṟṟēṣanṟe pradhāna kavāṭattil namukkŭ kāṇāṁ.}
\end{quote}

\begin{quote}
Let's meet at the station's main entrance.
\end{quote}

A formal word, used for big gates: a temple, a station, a park.

\begin{cousinweb}
\textbf{\ml{കവാടം}} (\emph{kavāṭaṁ}) is borrowed from Sanskrit \emph{kavāṭa}, \textquotedblleft{}the leaf of a door, a door\textquotedblright{} (Monier-Williams).
\end{cousinweb}

\begin{grammarlens}[title={What you've built}]
One more word for getting around.
\end{grammarlens}

\subsection*{Your turn}
\begin{itemize}
\raggedright
  \item \emph{Say it:} \emph{kavāṭaṁ}
  \item \emph{Say it:} \emph{kavāṭaṁ}, once more
  \item \emph{Say it:} \emph{sṟṟēṣanṟe pradhāna kavāṭattil namukkŭ kāṇāṁ.} — Let's meet at the station's main entrance.
  \item \emph{Recall:} say the Malayalam for a lift, then the Malayalam for a staircase, then say \emph{kavāṭaṁ} again
\end{itemize}

\begin{tcolorbox}[breakable,colback=teal!4,colframe=teal!35!black,title={Before you move on}]
What does \ml{കവാടം} mean? (\textquotedblleft{}A gateway; an entrance\textquotedblright{}.) Say it once more.
\end{tcolorbox}
\section[\texorpdfstring{\ml{വിശ്രമമുറി} (viśramamuṟi)}{വിശ്രമമുറി (viśramamuṟi)}]{\ml{വിശ്രമമുറി} (viśramamuṟi) — a waiting room; a rest room}

\label{lesson:ML-C485-visramamuri}



\begin{quote}
Before the new one: say the Malayalam for a staircase, then the Malayalam for a gateway.
\end{quote}

\subsection*{You'll want to know: \ml{വിശ്രമമുറി}}
\textbf{\ml{വിശ്രമമുറി}} — \emph{viśramamuṟi} — \textquotedblleft{}a waiting room; a rest room\textquotedblright{}.

\begin{quote}
\ml{തീവണ്ടി} \ml{വൈകിയതുകൊണ്ട്} \ml{ഞങ്ങൾ} \ml{വിശ്രമമുറിയിൽ} \ml{ഇരുന്നു}.
\end{quote}

\begin{quote}
\emph{tīvaṇṭi vaikiyatukoṇṭŭ ñaṅṅaḷ viśramamuṟiyil irunnu.}
\end{quote}

\begin{quote}
Because the train was late, we sat in the waiting room.
\end{quote}

\ml{വിശ്രമം}, rest, plus \ml{മുറി}, a room.

\begin{cousinweb}
\textbf{\ml{വിശ്രമമുറി}} (\emph{viśramamuṟi}) is a Malayalam compound built from \ml{വിശ്രമം}, rest, from Sanskrit \emph{viśrama} (CDIAL 11947), and the native word \ml{മുറി}, a room (DEDR 5008).
\end{cousinweb}

\begin{grammarlens}[title={What you've built}]
One more word for getting around.
\end{grammarlens}

\subsection*{Your turn}
\begin{itemize}
\raggedright
  \item \emph{Say it:} \emph{viśramamuṟi}
  \item \emph{Say it:} \emph{viśramamuṟi}, once more
  \item \emph{Say it:} \emph{tīvaṇṭi vaikiyatukoṇṭŭ ñaṅṅaḷ viśramamuṟiyil irunnu.} — Because the train was late, we sat in the waiting room.
  \item \emph{Recall:} say the Malayalam for a staircase, then the Malayalam for a gateway, then say \emph{viśramamuṟi} again
\end{itemize}

\begin{tcolorbox}[breakable,colback=teal!4,colframe=teal!35!black,title={Before you move on}]
What does \ml{വിശ്രമമുറി} mean? (\textquotedblleft{}A waiting room; a rest room\textquotedblright{}.) Say it once more.
\end{tcolorbox}
\section[\texorpdfstring{\ml{തട്ടുകട} (taṭṭukaṭa)}{തട്ടുകട (taṭṭukaṭa)}]{\ml{തട്ടുകട} (taṭṭukaṭa) — a roadside food stall}

\label{lesson:ML-C485-tattukata}



\begin{quote}
Before the new one: say the Malayalam for a gateway, then the Malayalam for a waiting room.
\end{quote}

\subsection*{You'll want to know: \ml{തട്ടുകട}}
\textbf{\ml{തട്ടുകട}} — \emph{taṭṭukaṭa} — \textquotedblleft{}a roadside food stall\textquotedblright{}.

\begin{quote}
\ml{രാത്രി} \ml{ഞങ്ങൾ} \ml{ഒരു} \ml{തട്ടുകടയിൽ} \ml{നിന്ന്} \ml{ദോശ} \ml{കഴിച്ചു}.
\end{quote}

\begin{quote}
\emph{rātri ñaṅṅaḷ oru taṭṭukaṭayil ninnŭ dōśa kaḻiccu.}
\end{quote}

\begin{quote}
At night we ate dosa at a roadside stall.
\end{quote}

A small stall that cooks at night by the road, with a hot plate and a bench.

\begin{cousinweb}
\textbf{\ml{തട്ടുകട}} (\emph{taṭṭukaṭa}) is a Malayalam compound built from the native words \ml{തട്ട്}, a flat tray or platform (DEDR 3035), and \ml{കട}, a shop (DEDR 1142).
\end{cousinweb}

\begin{grammarlens}[title={What you've built}]
One more word for getting around.
\end{grammarlens}

\subsection*{Your turn}
\begin{itemize}
\raggedright
  \item \emph{Say it:} \emph{taṭṭukaṭa}
  \item \emph{Say it:} \emph{taṭṭukaṭa}, once more
  \item \emph{Say it:} \emph{rātri ñaṅṅaḷ oru taṭṭukaṭayil ninnŭ dōśa kaḻiccu.} — At night we ate dosa at a roadside stall.
  \item \emph{Recall:} say the Malayalam for a gateway, then the Malayalam for a waiting room, then say \emph{taṭṭukaṭa} again
\end{itemize}

\begin{tcolorbox}[breakable,colback=teal!4,colframe=teal!35!black,title={Before you move on}]
What does \ml{തട്ടുകട} mean? (\textquotedblleft{}A roadside food stall\textquotedblright{}.) Say it once more.
\end{tcolorbox}
